\section{The major contribution and terminal budget}
\label{sec:major-budget}

We now estimate the Taylor major and the full-variance tail, then assemble the regime-dependent sector partition.  The order of choices is essential: the Gaussian cutoff $C$ is fixed first, and only then is $X$ taken sufficiently large.  All constants are fixed once $t$, $\gamma$, $b$ and the stated numerical choices are fixed.

\subsection{A uniform local expansion}

Compactness of the Bernoulli interval $I_{t,\gamma}$ gives constants $\rho>0$ and $M<\infty$ such that, uniformly for $\theta\in I_{t,\gamma}$ and $|z|\leq\rho$,
\begin{equation}
\label{eq:local-log-expansion}
 \log\bigl((1-\theta)+\theta\exp(2\pi iz)\bigr)
 =2\pi i\theta z
 -2\pi^2\theta(1-\theta)z^2
 +R_\theta(z),
\end{equation}
where
\begin{equation}
\label{eq:local-log-remainder}
 |R_\theta(z)|\leq M|z|^3.
\end{equation}

Let
\[
 e_{\min}=\min_{e\in\mathcal E}e\asymp\min(X^2,Z).
\]
For $|m|\leq N=C/\sigma_{\mathcal E}$,
\begin{equation}
\label{eq:aggregate-cubic-general}
 \sum_{e\in\mathcal E}|m/e|^3
 \leq\frac{|m|^3}{e_{\min}}
       \sum_{e\in\mathcal E}\frac1{e^2}
 \ll\frac{C^3}{\sigma_{\mathcal E}e_{\min}}.
\end{equation}
By \cref{prop:variance-regimes}, this is
\begin{equation}
\label{eq:aggregate-cubic-remainder}
 \begin{cases}
 O(C^3\log X/X),&\gamma>2,\\
 O(C^3\sqrt{\log Z/Z}),&1<\gamma\leq2,
 \end{cases}
\end{equation}
and hence tends to zero for fixed $C$ in every regime.  Also $N=o(e_{\min})$, so all major phases lie in the common Taylor disk.

\subsection{Sector I: positive real major}

For $|m|\leq N$, \cref{eq:N-below-full} and the decoder-identification propositions show that the decoded tuple is the genuine integer frequency $m$ modulo $L$.  Summing \cref{eq:local-log-expansion} over the actual family, the linear term cancels exactly:
\[
 2\pi i m\theta\Lambda-2\pi i mt=0.
\]
The numerator $a$ has therefore introduced no residual phase.  The quadratic term is exactly the full variance.  Uniformly for $|m|\leq N$,
\begin{equation}
\label{eq:major-local-gaussian}
 F(m)=\exp\bigl(-2\pi^2m^2\sigma_{\mathcal E}^2+\varepsilon_m\bigr),
 \qquad
 \max_{|m|\leq N}|\varepsilon_m|\longrightarrow0.
\end{equation}
Thus, for all sufficiently large $X$,
\[
 \RePart F(m)\geq c_0\exp(-2\pi^2m^2\sigma_{\mathcal E}^2)
\]
with $c_0>0$ fixed.  Summing on $|m|\leq1/(4\sigma_{\mathcal E})$ gives:

\begin{proposition}[Positive Gaussian major]
\label{prop:positive-major}
For every fixed $C\geq1$ and all sufficiently large $X$,
\begin{equation}
\label{eq:positive-major}
 \RePart\sum_{|m|\leq N}F(m)
 \geq\frac{c_{\mathrm{maj}}}{\sigma_{\mathcal E}},
\end{equation}
where $c_{\mathrm{maj}}>0$ is independent of $C$ and $X$.
\end{proposition}

\subsection{Sector II: the full actual-variance tail}

For
\[
 N<|m|\leq T_{\mathrm{full}},
\]
the decoded tuple is the genuine integer $m$ modulo $L$.  The constant in \cref{eq:T-full} was chosen so that $|m|/e\leq1/4$ for every $e\in\mathcal E$.  Hence
\begin{equation}
\label{eq:full-variance-gaussian}
 |F(m)|\leq\exp(-c m^2\sigma_{\mathcal E}^2).
\end{equation}
A Gaussian integral comparison gives
\begin{equation}
\label{eq:sector-II-tail}
 \sum_{|m|>N}\exp(-c m^2\sigma_{\mathcal E}^2)
 \leq K\sigma_{\mathcal E}^{-1}\exp(-cC^2/2),
\end{equation}
where $K$ is independent of $C$ and $X$.  Choose $C$ so large that this is at most $c_{\mathrm{maj}}/(8\sigma_{\mathcal E})$.

\subsection{Strict assembly in all regimes}

Fix this $C$, then choose $X$ above the finite maximum of the thresholds required for:
\begin{itemize}
\item the parameterized prime counts, reciprocal masses, family size and complete square load;
\item separation from the support of $b$ and avoidance of $\mathcal T$;
\item the strict load inequalities $t<\Lambda<1$ and the compact Bernoulli interval;
\item anchor rigidity, fingerprint entropy and the weighted anchor partition;
\item both row-distance estimates, the row-tail sum and the two identification ranges;
\item the prime-only transition estimate when $\gamma\leq2$;
\item containment and prime supply in the adaptive interval;
\item the Taylor disk and the regime-specific cubic remainder;
\item all $o(\sigma_{\mathcal E}^{-1})$ totals in \cref{eq:transition-total,eq:adaptive-total,eq:anchor-tail-total,eq:terminal-error-total}.
\end{itemize}
The terminal order is therefore
\begin{equation}
\label{eq:parameter-order}
 (t,\gamma,b,\mathcal T)\ \text{fixed},
 \qquad C\longrightarrow X.
\end{equation}
For $\gamma>2$, Sectors III--V contribute $o(\sigma_{\mathcal E}^{-1})$.  For $1<\gamma\leq2$, Sectors III--VI do so, including the prime-only transition.  Sector II is at most $c_{\mathrm{maj}}/(8\sigma_{\mathcal E})$, whereas Sector I has real part at least $c_{\mathrm{maj}}/\sigma_{\mathcal E}$.  Applying \cref{prop:decoded-skeleton-positivity} yields:

\begin{theorem}[Positive parameterized Fourier numerator]
\label{thm:positive-fourier-numerator}
For every fixed reduced $t=a/b\in(0,1)$ with squarefree $b$, every fixed $\gamma$ satisfying \cref{eq:admissible-region}, and every finite forbidden set $\mathcal T$, one has for all sufficiently large $X$
\begin{equation}
\label{eq:positive-full-sum}
 \RePart\sum_{h\bmod L}F(h)>0.
\end{equation}
\end{theorem}

The proof uses the full actual-family variance in both the major and the first minor range.  At $\gamma=2$, that variance is target-row dominated, while the following transition range is controlled solely by the complete prime-pair square load.  No estimate is obtained by formal substitution for the exponent $3$.
